\documentclass[11pt,a4paper]{article}  % Déclare la classe du document.

\newcommand{\chaptermark}[2]{\og #2 \fg{} (#1)} % ligne uniquement pour éviter une erreur en cas d'article
\include{../1ere-preambule}
\usepackage{mathrsfs}

\rhead{Exercices d'entrainement - Application du produit scalaire}

\newcommand{\va}[1]{\lvert \ #1 \ \rvert}
\newcommand{\vaf}[1]{\left\lvert \ #1\  \right\rvert}

\begin{document}
	\graphicspath{{../images/}}
	
	\begin{center}
		
		
		\fbox{
			\begin{minipage}[]{15cm}
				\begin{center}
					\LARGE{\rule{0.5cm}{0pt}\rule[-0.25cm]{0pt}{0.9cm}
						Exercices d'entrainement\\
						Application du produit scalaire
					}
				\end{center}
			\end{minipage}
		}
	\end{center}

\exerciceDS{}
$\mathrm{ABC}$ est un triangle tel que :$A B=2, A C=6 \text { et } \widehat{B A C}=60^{\circ} .$\\
Calculer la valeur exacte de la longueur $B C$.


\exerciceDS{}
ABCD est un parallélogramme tel que :$A B=6, A D=4 \text { et } \widehat{D A B}=60^{\circ} .$
\begin{multicols}{2}
	\begin{enumMet}
		\item Construire une figure.
		\item Calculer la longueur BD.
		\item Quelle est la mesure de l'angle $\widehat{A B C}$ ?
		\item Calculer la longueur AC.
	\end{enumMet}
\end{multicols}

\vspace{-0.3cm}
\exerciceDS{}

\begin{minipage}[c]{0.55\linewidth}
	Un voilier V approche de la côte. Il est situé à 9 milles du phare A, à 5 milles du phare $B$ et l'angle $\widehat{AVB}$ mesure $60^{\circ} .$
	Calculer la distance AB entre les deux phares. Arrondir au dixième.
\end{minipage} \hfill
\begin{minipage}[c]{0.4\linewidth}
	\begin{center}
		\includegraphics[scale=0.9]{../images/1.06.04 - Img01.png}
	\end{center}
\end{minipage}

\exerciceDS{}
Dans les 2 cas ci dessous, calculer $\vv{A B} \cdot \vv{A C}$.
\begin{multicols}{2}
	\begin{enumMet}
		\item \ \\
		\includegraphics[scale=1.2]{../images/1.06.04 - Img02.png}
		\item \ \\
		\includegraphics[scale=1.]{../images/1.06.04 - Img03.png}
	\end{enumMet}
\end{multicols}

\exerciceDS{}
On considère deux vecteurs $\vv{u}$ et $\vv{v}$ du plan tels que $\|\vv{u}+\vv{v}\|=5$ et $\|\vv{u}-\vv{v}\|=3$. Calculer le produit scalaire $\vv{u} \cdot \vv{v}$.

\exerciceDS{}
Dans le plan muni d'un repère orthonormé, on considère les vecteurs
$\Coord{u}{2}{-3}$ et $\Coord{v}{1}{2}$.
\begin{enumMet}
	\item calculer $\vv{u} \cdot \vv{v},\|\vv{u}\|$ et $\|\vv{v}\|$.
	\item Calculer de deux façons les nombres : $\|\vv{u}+\vv{v}\|^{2} \text { et }\|\vv{u}-\vv{v}\|^{2}$
	\item Calculer de deux façons $(\vv{u}+\vv{v}) \cdot(\vv{u}-\vv{v})$.
\end{enumMet}

\exerciceDS{}
Le plan est muni d'un repère orthonormé. Dans chaque cas, calculer $\vv{A B} \cdot \vv{A C}$ puis une mesure (arrondie au degré près) de l'angle $\widehat{B A C}$.
\begin{enumMet}
	\item $A(-1 ; 4), B(2 ; 2)$ et $C(1 ; 5)$
	\item $A(2 ; 4), B(1 ; 3)$ et $C(4 ; 2)$
	\item $A(0 ; 4), B(-3 ; 5)$ et $C(6 ; 2)$
\end{enumMet}


\exerciceDS{}
Dans un repère orthonormé du plan, on considère les points $A(1 ; 4), B(-3 ;-2)$ et $C(2 ;-5)$.\\
Déterminer les angles du triangle $A B C$ à 0,1 degré près.

\exerciceDS{}
$A B C$ est un triangle tel que $A B=4, A C=3$ et $\widehat{B A C}=135^{\circ} . M$ est le milieu du segment $[A C]$ et $P$ est le point déini par $\vv{A P}=\frac{3}{4} \vv{A B}$.
\begin{enumMet}
	\item En notant que $\vv{M P}=\vv{A P}-\vv{A M}$, calculer $\vv{M P^{2}}$.
	\item En déduire la longueur $M P$.
\end{enumMet}

\exerciceDS{}
$A B C D$ est un parallélogramme tel que $A B=3$, $A D=5$ et $A C=6$.
\begin{enumMet}
	\item Rappeler la formule exprimant $\vv{u} \cdot \vv{v}$ en fonction de $\|\vv{u}\|,\|\vv{v}\|$ et $\|\vv{u}+\vv{v}\|$.\\
	Calculer alors le produit scalaire $\vv{A B} \cdot \vv{A D}$.
	\item En déduire une mesure, en degré, de l'angle $\widehat{B A D}$ (on donnera une valeur arrondie au degré près).
\end{enumMet}

\exerciceDS{}
$A B C D$ est un parallélogramme tel que $A B=4$, $A D=6$ et $\widehat{B A D}=40^{\circ}$.\\
Calculer le produit scalaire $\vv{A B} \cdot \vv{A C}$ (on pourra décomposer le vecteur $\vv{A C}$ ).

\exerciceDS{}

Donner un vecteur directeur pour chacune des droites suivantes et en déduire qu'elles sont perpendiculaires.
\begin{enumMet}
	\item Pour les droites $d_1$ et $d_2$ d'équations cartésiennes $2 x-3 y+4=0$ et $3 x+2 y-1=0$
	\item Pour les droites $d_1$ et $d_2$ d'équations cartésiennes $x-y+3=0$ et $2 x+2 y-1=0$.
	\item Pour les droites $d_1$ et $d_2$ d'équations $y=-3 x+1$ et $-x+3 y-1=0$
\end{enumMet}



\exerciceDS{}

Déterminer un vecteur normal à chacune des droites données par les équations cartésiennes ci-dessous.
\begin{enumMet}
	\item $2 x+y-3=0$
	\item $-3 x+5 y=0$
	\item $5 x-3 y+2=0$
	\item $-2 x-4 y+1=0$
\end{enumMet}

\exerciceDS{}

Déterminer un vecteur normal à chacune des droites données par les équations réduites ci-dessous.
\begin{enumMet}
	\item $y=2 x-3$
	\item $y=4 x$
	\item $y=-3 x+1$
	\item $y=-7 x-2$
\end{enumMet}

\exerciceDS{}

Déterminer un vecteur normal à chacune des droites définies par les deux points donnés.
\begin{enumMet}
	\item $\mathrm{B}(-3 ; 2)$ et $\mathrm{C}(1 ;-2)$
	\item $\mathrm{F}(1 ; 0)$ et $\mathrm{G}(-3 ; 4)$
	\item $\mathrm{M}(0 ;-2)$ et $\mathrm{N}(5 ; 4)$
	\item $\mathrm{H}(-2 ; 3)$ et $\mathrm{K}(-1 ;-5)$
\end{enumMet}

\exerciceDS{}

Dans chacun des cas, donner une équation cartésienne de la droite passant par le point $\mathrm{B}$ et de vecteur normal $\vec{n}$.
\begin{enumMet}
	\item $B(4 ; 1)$ et $\Coord{n}{1}{-3}$
	\item $\mathrm{B}(2 ; 3)$ et $\Coord{n}{4}{2}$
	\item $\mathrm{B}(-2 ;-1)$ et $\Coord{n}{-3}{-1}$
	\item $\mathrm{B}(0 ; 5)$ et $\Coord{n}{5}{3}$
\end{enumMet}

\exerciceDS{}
On considère la droite $d$ d'équation réduite $y=-4 x+5$.
\begin{enumMet}
	\item Donner un vecteur directeur de $d$.
	\item En déduire un vecteur normal de $d$.
	\item Donner une équation cartésienne de la droite $d$ ' perpendiculaire à $d$ passant par le point $\mathrm{A}(1 ; 1)$.
	\item Donner une équation réduite de la droite $d^{\prime}$.
\end{enumMet}

\exerciceDS{}
On considère la droite $d$ d'équation cartésienne $-5 x+3 y-1=0$.
\begin{enumMet}
	\item 1. Donner un vecteur directeur de $d$.
	\item En déduire un vecteur normal de $d$.
	\item Donner une équation cartésienne de la droite $d^{\prime}$ perpendiculaire à $d$ passant par le point $\mathrm{F}(2 ;-2)$.
	\item Donner une équation réduite de la droite $d^{\prime}$.
\end{enumMet}


\exerciceDS{} %Ex 81
On donne deux points $A(-1 ; 2)$ et $\mathrm{B}(3 ; 4)$.
\begin{enumMet}
	\item Déterminer les coordonnées du milieu I du segment [AB].
	\item Déterminer les coordonnées du vecteur $\vv{A B}$.
	\item Déterminer une équation cartésienne de la médiatrice du segment $[A B]$.
\end{enumMet}


\exerciceDS{} %Ex 85
On considère un triangle $A B C$ avec $A(1 ;-3), B(-2 ; 4)$ et $\mathrm{C}(3 ; 0)$.
\begin{enumMet}
	\item 1. Déterminer les coordonnées du vecteur $\vv{B C}$.
	\item Déterminer une équation cartésienne de la hauteur issue de $A$ dans le triangle $A B C$.
	\item Déterminer les coordonnées de $\mathrm{H}$ projeté orthogonal de A sur (BC).
	\item Calculer la longueur $\mathrm{AH}$.
	\item En déduire l'aire du triangle $A B C$.
\end{enumMet}

\exerciceDS{} %Ex 86
On donne les points $E(-2 ;-2), F(5 ;-4)$ et $\mathrm{G}(1 ; 2)$.
\begin{enumMet}
	\item Déterminer une équation cartésienne de la médiatrice du segment [EF].
	\item Le point $\mathrm{G}$ appartient-il à cette droite ?
	\item Déterminer une équation de la hauteur issue de $\mathrm{G}$ dans le triangle EFG.
	\item Déterminer les coordonnées de K projeté orthogonal de $G$ sur la droite (EF).
	\item Calculer l'aire du triangle EFG.
\end{enumMet}


\end{document}

